\section*{ID9211: Definition of s1(t,Ω)}
\paragraph{Metadata.}
PiID: 8892186543648 | RTEID: RTE:P35768.3586:T2.3866 | UUID: a7a5dba7-aad8-5be4-9882-615d6ec38e9b

\paragraph{Canonical Statement.}
dancy. (3) Enables projection onto toroidal fibers via Canonical Registry Vol. 01 103 Trawin Zero Point Contextus Canonicus 2 TZP CANONICAL REGISTRY – ENTRIES 150–199 consistent opposite mappings. Falsifiability / Test Handle: Enumerate all addresses and verify that each has a unique opposite by modular arithmetic; confirm that factorization yields integer refinement tiers. 2.0.154 ID 303: Rotating Double Helix Geometry Type: Definition Formal Statement: Two helical strands are parametrized by \$s1(t,Ω)=R(Ωt)(Rcos(2πt/P), Rsin(2πt/P),\$ t)s\_1(t,\textbackslash{}Omega)=R(\textbackslash{}Omega t)(R \textbackslash{}cos(2\textbackslash{}pi t/P),\textbackslash{} R \textbackslash{}sin(2\textbackslash{}pi t/P),\textbackslash{} \$t)s1(t,Ω)=R(Ωt)(Rcos(2πt/P), Rsin(2πt/P),\$ t) and \$s2(t,Ω)=R(Ωt)(Rcos(2πt/P+π), Rsin(2πt/P+π),\$ t)s\_2(t,\textbackslash{}Omega)=R(\textbackslash{}Omega t)(R \textbackslash{}cos(2\textbackslash{}pi t/P + \textbackslash{}pi),\textbackslash{} R \textbackslash{}sin(2\textbackslash{}pi t/P + \textbackslash{}pi),\textbackslash{} \$t)s2(t,Ω)=R(Ωt)(Rcos(2πt/P+π), Rsin(2πt/P+π),\$ t), where \$R(θ)R(\textbackslash{}theta)R(θ)\$ is the rotation matrix about the zzz-axis, Ω\textbackslash{}OmegaΩ the angular rotation, RR

\paragraph{Canonical Equation.}
\[
s1(t,Ω)=R(Ωt)(Rcos(2πt/P), Rsin(2πt/P),
\]

\paragraph{Dictionary.}
Definition of s1(t,Ω) — s1(t,Ω)=R(Ωt)(Rcos(2πt/P), Rsin(2πt/P)

\paragraph{Encyclopedia.}
Definition of s1(t,Ω) is an algebraic definition, written s1(t,Ω)=R(Ωt)(Rcos(2πt/P), Rsin(2πt/P). It is an algebraic definition expressing the quantity directly in terms of the variables on its right-hand side. It is built from R, c, o, s, P, n, defining s1(t,Ω). Within the Trawin Zero Point registry it serves as one element of the framework's mathematical narrative.
